\documentclass[10pt,a4paper]{amsart}
\usepackage{amsmath,amssymb,amsthm,graphicx}
\usepackage{verbatim}
\newenvironment{reference}{\comment}{\endcomment}
%\newenvironment{reference}{}{}
\newenvironment{slogan}{\comment}{\endcomment}
\newenvironment{history}{\comment}{\endcomment}

% For commutative diagrams we use Xy-pic
\usepackage[all]{xy}

% We use 2cell for 2-commutative diagrams.
\xyoption{2cell}
% Theorem environments.
%
\theoremstyle{plain}
\newtheorem{theorem}[subsection]{Theorem}
\newtheorem{proposition}[subsection]{Proposition}
\newtheorem{lemma}[subsection]{Lemma}

\theoremstyle{definition}
\newtheorem{definition}[subsection]{Definition}
\newtheorem{example}[subsection]{Example}
\newtheorem{exercise}[subsection]{Exercise}
\newtheorem{situation}[subsection]{Situation}

\theoremstyle{remark}
\newtheorem{remark}[subsection]{Remark}
\newtheorem{remarks}[subsection]{Remarks}

\numberwithin{equation}{subsection}

% Macros
%
\def\lim{\mathop{\mathrm{lim}}\nolimits}
\def\colim{\mathop{\mathrm{colim}}\nolimits}
\def\Spec{\mathop{\mathrm{Spec}}}
\def\Hom{\mathop{\mathrm{Hom}}\nolimits}
\def\Ext{\mathop{\mathrm{Ext}}\nolimits}
\def\SheafHom{\mathop{\mathcal{H}\!\mathit{om}}\nolimits}
\def\SheafExt{\mathop{\mathcal{E}\!\mathit{xt}}\nolimits}
\def\Sch{\mathit{Sch}}
\def\Mor{\mathop{\mathrm{Mor}}\nolimits}
\def\Ob{\mathop{\mathrm{Ob}}\nolimits}
\def\Sh{\mathop{\mathit{Sh}}\nolimits}
\def\NL{\mathop{N\!L}\nolimits}
\def\CH{\mathop{\mathrm{CH}}\nolimits}
\def\proetale{{pro\text{-}\acute{e}tale}}
\def\etale{{\acute{e}tale}}
\def\QCoh{\mathit{QCoh}}
\def\Ker{\mathop{\mathrm{Ker}}}
\def\Im{\mathop{\mathrm{Im}}}
\def\Coker{\mathop{\mathrm{Coker}}}
\def\Coim{\mathop{\mathrm{Coim}}}

% Boxtimes
%
\DeclareMathSymbol{\boxtimes}{\mathbin}{AMSa}{"02}

%
% Macros for moduli stacks/spaces
%
\def\QCohstack{\mathcal{QC}\!\mathit{oh}}
\def\Cohstack{\mathcal{C}\!\mathit{oh}}
\def\Spacesstack{\mathcal{S}\!\mathit{paces}}
\def\Quotfunctor{\mathrm{Quot}}
\def\Hilbfunctor{\mathrm{Hilb}}
\def\Curvesstack{\mathcal{C}\!\mathit{urves}}
\def\Polarizedstack{\mathcal{P}\!\mathit{olarized}}
\def\Complexesstack{\mathcal{C}\!\mathit{omplexes}}
% \Pic is the operator that assigns to X its picard group, usage \Pic(X)
% \Picardstack_{X/B} denotes the Picard stack of X over B
% \Picardfunctor_{X/B} denotes the Picard functor of X over B
\def\Pic{\mathop{\mathrm{Pic}}\nolimits}
\def\Picardstack{\mathcal{P}\!\mathit{ic}}
\def\Picardfunctor{\mathrm{Pic}}
\def\Deformationcategory{\mathcal{D}\!\mathit{ef}}
\usepackage{cleveref}
\crefname{originalreference}{original reference}{original references}
\Crefname{originalreference}{Original reference}{Original references}
\makeatletter
\newlabel{situation-setup}{{02QL}{source}{}{sourcepointer-1}{}}
\newlabel{situation-setup@cref}{{[originalreference][0][]02QL}{[1][0][]source}{}{}{}}
\makeatother
\crefname{originalreference}{original reference}{original references}
\Crefname{originalreference}{Original reference}{Original references}
\makeatletter
\newlabel{remark-bivariant-commute}{{0FA2}{source}{}{sourcepointer-2}{}}
\newlabel{remark-bivariant-commute@cref}{{[originalreference][0][]0FA2}{[1][0][]source}{}{}{}}
\makeatother
\crefname{originallemma}{original lemma}{original lemmas}
\Crefname{originallemma}{Original lemma}{Original lemmas}
\makeatletter
\newlabel{lemma-bivariant-zero}{{02UC}{source}{}{sourcepointer-3}{}}
\newlabel{lemma-bivariant-zero@cref}{{[originallemma][0][]02UC}{[1][0][]source}{}{}{}}
\makeatother
\crefname{originallemma}{original lemma}{original lemmas}
\Crefname{originallemma}{Original lemma}{Original lemmas}
\makeatletter
\newlabel{lemma-gysin-fundamental}{{0FBM}{source}{}{sourcepointer-4}{}}
\newlabel{lemma-gysin-fundamental@cref}{{[originallemma][0][]0FBM}{[1][0][]source}{}{}{}}
\makeatother
\crefname{originallemma}{original lemma}{original lemmas}
\Crefname{originallemma}{Original lemma}{Original lemmas}
\makeatletter
\newlabel{lemma-cap-commutative-chern}{{02UA}{source}{}{sourcepointer-5}{}}
\newlabel{lemma-cap-commutative-chern@cref}{{[originallemma][0][]02UA}{[1][0][]source}{}{}{}}
\makeatother
\crefname{originallemma}{original lemma}{original lemmas}
\Crefname{originallemma}{Original lemma}{Original lemmas}
\makeatletter
\newlabel{divisors-lemma-blowing-up-gives-effective-Cartier-divisor}{{02OS}{source}{}{sourcepointer-6}{}}
\newlabel{divisors-lemma-blowing-up-gives-effective-Cartier-divisor@cref}{{[originallemma][0][]02OS}{[1][0][]source}{}{}{}}
\makeatother
\crefname{originallemma}{original lemma}{original lemmas}
\Crefname{originallemma}{Original lemma}{Original lemmas}
\makeatletter
\newlabel{lemma-gysin-bivariant}{{0B79}{source}{}{sourcepointer-7}{}}
\newlabel{lemma-gysin-bivariant@cref}{{[originallemma][0][]0B79}{[1][0][]source}{}{}{}}
\makeatother
\crefname{originallemma}{original lemma}{original lemmas}
\Crefname{originallemma}{Original lemma}{Original lemmas}
\makeatletter
\newlabel{lemma-gysin-decompose}{{0FBL}{source}{}{sourcepointer-8}{}}
\newlabel{lemma-gysin-decompose@cref}{{[originallemma][0][]0FBL}{[1][0][]source}{}{}{}}
\makeatother
\crefname{originallemma}{original lemma}{original lemmas}
\Crefname{originallemma}{Original lemma}{Original lemmas}
\makeatletter
\newlabel{lemma-gysin-agrees}{{0FBN}{source}{}{sourcepointer-9}{}}
\newlabel{lemma-gysin-agrees@cref}{{[originallemma][0][]0FBN}{[1][0][]source}{}{}{}}
\makeatother
\title{Every bivariant class commutes with the class of a virtual normal sheaf}
\author{The Stacks Project Authors}
\date{stacks-project:chow.tex:local-snapshot-20261012; GNU FDL 1.2 or later}
\begin{document}
\maketitle
\noindent\textit{Editorial note.} All delivered mathematical text is the author's own, copied from the pinned Stacks Project chapter file chow.tex; no statement or proof step was rewritten, reordered or removed.\par\medskip
\noindent\textit{Editorial note.} Preamble provenance: the pinned chapter file opens with an input of the shared preamble and carries no preamble text of its own, so the project's single Stacks-wide preamble staged at sources/stacks-chow/preamble.tex is excerpted verbatim as two spans, lines 14-18 (Xy-pic setup) and lines 154-217 (the native theorem declarations and mathematical macros). Nothing was added, removed or redefined; the wrapper class and the amsmath/amssymb/amsthm/graphicx base packages are supplied by the assembler.\par\medskip
\noindent\textit{Editorial note.} The counted statement and proof contain no cross reference, no citation and no included file, so the delivered body is self contained.\par\medskip
\noindent\textit{Editorial note.} Canonical location of the counted proposition: Stacks tag 0FBP (sources/stacks-chow/pinned-original-source.tex line 12206).\par\medskip
\noindent\textit{Allowed standard prerequisites (editorial).}\begin{itemize}
\item Testing identities of bivariant classes on fundamental classes. Same chapter, Lemma labelled lemma-bivariant-zero, cited at proof lines 12214-12216 and 12230-12232. Reduces to an integral X and allows the blowup reduction.
\item The fundamental computation identifying the class with a top Chern class in the equal case. Same chapter, Lemma labelled lemma-gysin-fundamental, cited at proof lines 12223-12226. Handles the case Z = X.
\item Decomposition of the class of a virtual normal sheaf along an effective Cartier divisor. Same chapter, Lemmas labelled lemma-gysin-decompose and lemma-gysin-agrees, cited at proof lines 12240-12245, and Lemma labelled lemma-gysin-bivariant for the Gysin homomorphism i\textasciicircum{}*. Writes the class as c\_top(E) o i\textasciicircum{}* in the divisor case.
\item Centrality of Chern classes among bivariant classes. Same chapter, Lemma labelled lemma-cap-commutative-chern, cited at proof lines 12226 and 12246-12248. Supplies the commutation with the top Chern class and with the product expression.
\end{itemize}\medskip
\ifdefined\hypertarget\hypertarget{sourcepointer-1}{}\fi
\noindent\textit{Source reference.} Original reference 02QL: sources/stacks-chow/pinned-original-source.tex, source lines 1372--1372. Same-chapter pointer into the author environment situation at source line 1372; invoked at proof line None of the counted proof. It is retained only as a navigation pointer; every load-bearing step is inside the delivered span.\par\medskip
\ifdefined\hypertarget\hypertarget{sourcepointer-2}{}\fi
\noindent\textit{Source reference.} Original reference 0FA2: sources/stacks-chow/pinned-original-source.tex, source lines 6200--6200. Same-chapter pointer into the author environment remark at source line 6200; invoked at proof line None of the counted proof. It is retained only as a navigation pointer; every load-bearing step is inside the delivered span.\par\medskip
\ifdefined\hypertarget\hypertarget{sourcepointer-3}{}\fi
\noindent\textit{Source reference.} Original lemma 02UC: sources/stacks-chow/pinned-original-source.tex, source lines 6585--6585. Same-chapter pointer into the author environment lemma at source line 6585; invoked at proof line 12215 of the counted proof. It is retained only as a navigation pointer; every load-bearing step is inside the delivered span.\par\medskip
\ifdefined\hypertarget\hypertarget{sourcepointer-4}{}\fi
\noindent\textit{Source reference.} Original lemma 0FBM: sources/stacks-chow/pinned-original-source.tex, source lines 12076--12076. Same-chapter pointer into the author environment lemma at source line 12076; invoked at proof line 12222 of the counted proof. It is retained only as a navigation pointer; every load-bearing step is inside the delivered span.\par\medskip
\ifdefined\hypertarget\hypertarget{sourcepointer-5}{}\fi
\noindent\textit{Source reference.} Original lemma 02UA: sources/stacks-chow/pinned-original-source.tex, source lines 7581--7581. Same-chapter pointer into the author environment lemma at source line 7581; invoked at proof line 12223 of the counted proof. It is retained only as a navigation pointer; every load-bearing step is inside the delivered span.\par\medskip
\ifdefined\hypertarget\hypertarget{sourcepointer-6}{}\fi
\noindent\textit{Source reference.} Original lemma 02OS: sources/stacks-chow/tags.txt, source lines 3527--3527. Cross-chapter standard prerequisite located in the project tag list staged in the local snapshot; invoked at proof line 12231 of the counted proof. It is retained only as a navigation pointer; every load-bearing step is inside the delivered span.\par\medskip
\ifdefined\hypertarget\hypertarget{sourcepointer-7}{}\fi
\noindent\textit{Source reference.} Original lemma 0B79: sources/stacks-chow/pinned-original-source.tex, source lines 6149--6149. Same-chapter pointer into the author environment lemma at source line 6149; invoked at proof line 12240 of the counted proof. It is retained only as a navigation pointer; every load-bearing step is inside the delivered span.\par\medskip
\ifdefined\hypertarget\hypertarget{sourcepointer-8}{}\fi
\noindent\textit{Source reference.} Original lemma 0FBL: sources/stacks-chow/pinned-original-source.tex, source lines 11994--11994. Same-chapter pointer into the author environment lemma at source line 11994; invoked at proof line 12243 of the counted proof. It is retained only as a navigation pointer; every load-bearing step is inside the delivered span.\par\medskip
\ifdefined\hypertarget\hypertarget{sourcepointer-9}{}\fi
\noindent\textit{Source reference.} Original lemma 0FBN: sources/stacks-chow/pinned-original-source.tex, source lines 12174--12174. Same-chapter pointer into the author environment lemma at source line 12174; invoked at proof line 12243 of the counted proof. It is retained only as a navigation pointer; every load-bearing step is inside the delivered span.\par\medskip
\begin{lemma}
\label{lemma-gysin-commutes}
Let $(S, \delta)$ be as in Situation \ref{situation-setup}. Let $X$ be a scheme
locally of finite type over $S$. Let $Z \subset X$ be a closed subscheme
with virtual normal sheaf $\mathcal{N}$. Let $Y \to X$ be locally of
finite type and $c \in A^*(Y \to X)$. Then $c$ and $c(Z \to X, \mathcal{N})$
commute (Remark \ref{remark-bivariant-commute}).
\end{lemma}

\begin{proof}
To check this we may use Lemma \ref{lemma-bivariant-zero}.
Thus we may assume $X$ is an integral scheme and we have to show
$c \cap c(Z \to X, \mathcal{N}) \cap [X] =
c(Z \to X, \mathcal{N}) \cap c \cap [X]$ in $\CH_*(Z \times_X Y)$.

\medskip\noindent
If $Z = X$, then $c(Z \to X, \mathcal{N}) = c_{top}(\mathcal{N})$ by
Lemma \ref{lemma-gysin-fundamental} which commutes
with the bivariant class $c$, see Lemma \ref{lemma-cap-commutative-chern}.

\medskip\noindent
Assume that $Z$ is not equal to $X$. By Lemma \ref{lemma-bivariant-zero}
it even suffices to prove the result after blowing up $X$ (in a nonzero ideal).
Let us blowup $X$ in the ideal sheaf of $Z$. This reduces us to the case
where $Z$ is an effective Cartier divisor, see
Divisors, Lemma
\ref{divisors-lemma-blowing-up-gives-effective-Cartier-divisor},

\medskip\noindent
If $Z$ is an effective Cartier divisor, then we have
$$
c(Z \to X, \mathcal{N}) =
c_{top}(\mathcal{E}) \circ i^*
$$
where $i^* \in A^1(Z \to X)$ is the Gysin homomorphism
associated to $i : Z \to X$ (Lemma \ref{lemma-gysin-bivariant})
and $\mathcal{E}$ is the dual of the kernel of
$\mathcal{N}^\vee \to \mathcal{C}_{Z/X}$, see
Lemmas \ref{lemma-gysin-decompose} and \ref{lemma-gysin-agrees}.
Then we conclude because Chern classes are in the center of the
bivariant ring (in the strong sense formulated in
Lemma \ref{lemma-cap-commutative-chern}) and $c$ commutes
with the Gysin homomorphism $i^*$ by definition of bivariant classes.
\end{proof}

\end{document}
